% ============================================================================
%  Rozdział 5. Testy i walidacja
% ============================================================================
\chapter{Testy i walidacja}
\label{rozdz:testy}

Testy miały na celu sprawdzenie, czy gra spełnia wymagania określone w~rozdziale trzecim,
czy przepływ pomiędzy scenami i~łańcuchy zdarzeń działają poprawnie oraz czy rozgrywka
jest czytelna dla gracza. Zastosowano cztery uzupełniające się metody: testy funkcjonalne
w~trybie odtwarzania edytora Unity, automatyczny scenariusz uruchomieniowy w~trybie
wsadowym edytora, walidację konfiguracji w~procesie budowania oraz sesje testowe rozgrywki,
po których wprowadzano poprawki w~mechanikach.

\section{Środowisko i metoda testów}

Testy funkcjonalne i~sesje testowe rozgrywki przeprowadzono w~edytorze Unity 2022.3.60f1
na dwóch komputerach: stacjonarnym z~procesorem AMD Ryzen~5 9600X, 32~GB pamięci
operacyjnej i~kartą graficzną AMD Radeon RX~9070~XT (podrozdział~3.4.1) oraz na laptopie
HP, z~użyciem myszy jako urządzenia wejściowego. Uruchomienie gry na dwóch maszynach
o~wyraźnie różnej wydajności pozwoliło sprawdzić, że rozgrywka nie zależy od liczby klatek
obrazu, ponieważ ruch kuli i~przeszkód jest obliczany w~stałym kroku fizyki.
W~trakcie sesji testowych oceniano przede wszystkim wygodę sterowania, czytelność interfejsu
i~poziom trudności, a~ich wyniki opisano w~podrozdziale~\ref{sec:poprawki}.

\section{Testy funkcjonalne}

Scenariusze testowe przygotowano na podstawie wymagań funkcjonalnych. Każdy scenariusz
opisuje czynność gracza lub zdarzenie rozgrywki oraz oczekiwaną reakcję gry. Wyniki
zestawiono w~tabeli~\ref{tab:testy}. Scenariusze oznaczone literą A~zostały dodatkowo
wykonane automatycznie w~scenariuszu wsadowym opisanym w~kolejnym podrozdziale.

\begin{table}[!htb]
  \centering
  \caption{Wyniki testów funkcjonalnych; \zrodloWlasne}
  \label{tab:testy}
  \begin{tabularx}{\textwidth}{c >{\raggedright\arraybackslash}p{5.0cm} X c}
    \toprule
    \textbf{Nr} & \textbf{Scenariusz} & \textbf{Oczekiwany rezultat} & \textbf{Wynik} \\
    \midrule
    1 & Uruchomienie gry (A) & menu z~przyciskami trzech poziomów, statystyk i~wyjścia & zgodny \\
    2 & Otwarcie okna statystyk (A) & tabela trzech poziomów z~aktualnymi wartościami & zgodny \\
    3 & Wybór poziomu w~menu (A) & ekran ładowania, a~następnie wybrany poziom & zgodny \\
    4 & Rozpoczęcie poziomu (A) & kula w~punkcie startowym, licznik czasu odlicza, zdrowie 100 & zgodny \\
    5 & Gest rozpoczęty poza kulą & brak strzału i~linii celowniczej & zgodny \\
    6 & Gest rozpoczęty na kuli & linia celownicza, po zwolnieniu ruch w~kierunku przeciwnym & zgodny \\
    7 & Uderzenie kuli w~ścianę & odbicie, animacja zgniecenia, efekt pyłu & zgodny \\
    8 & Zebranie monety & animacja znikania, wzrost licznika o~20 & zgodny \\
    9 & Zebranie ostatniej monety & pojawienie się mety & zgodny \\
    10 & Dotarcie do mety (A) & zatrzymanie czasu, okno z~czasem i~informacją o~rekordzie & zgodny \\
    11 & Kolizja z~przeciwnikiem & odbicie z~utratą 80\% mocy, obrażenia zależne od typu przeciwnika & zgodny \\
    12 & Zadanie kuli 29 punktów obrażeń (A) & aktualizacja paska i~liczby punktów życia & zgodny \\
    13 & Wjazd w~apteczkę przy pełnym zdrowiu & apteczka pozostaje na planszy & zgodny \\
    14 & Utrata zdrowia (A) & kula znika, okno końca bez panelu czasu & zgodny \\
    15 & Wjazd na pole przyspieszające & wyrzucenie kuli wzdłuż strzałek & zgodny \\
    16 & Naciśnięcie \emph{Push} przy stojącej kuli & ruch w~losowym kierunku; przy toczącej się kuli przycisk nieaktywny & zgodny \\
    17 & Ponowna gra i~powrót do menu & wczytanie poziomu od początku lub menu, zapis czasu gry & zgodny \\
    \bottomrule
  \end{tabularx}
\end{table}

Szczególną uwagę poświęcono scenariuszom, w~których jedno zdarzenie uruchamia kilka
reakcji w~różnych modułach. Przy wjeździe na metę licznik czasu zapisuje wynik przed
dezaktywacją kuli, więc okno końca poziomu zawsze pokazuje czas przejścia
(rysunek~\ref{rys:wyniki}a). Przy utracie zdrowia wynik nie jest publikowany, dzięki
czemu to samo okno wyświetla się bez panelu czasu (rysunek~\ref{rys:obrazenia}b).

\section{Automatyczny scenariusz uruchomieniowy}

Aby sprawdzić przepływ scen bez udziału człowieka, przygotowano skrypt edytora
uruchamiany na kopii projektu w~trybie wsadowym (\kod{-batchmode}), z~wykorzystaniem
procesora graficznego. Skrypt otwiera menu główne, przechodzi w~tryb odtwarzania
i~wykonuje kolejne kroki: otwiera okno statystyk, wybiera pierwszy poziom tym samym
komponentem, którego używa przycisk w~menu, a~po wczytaniu poziomu nadaje kuli prędkość
operacją \kod{Launch} i~publikuje zdarzenie dotarcia do mety. Następnie wczytuje drugi
poziom, zadaje kuli dwukrotnie 29~punktów obrażeń przez komponent \klasa{BallGetDamage},
a~potem obrażenia przekraczające pozostały zapas zdrowia, i~na koniec wczytuje trzeci
poziom. Po każdym kroku skrypt renderuje obraz kamery głównej wraz z~interfejsem do pliku;
zrzuty ekranu zamieszczone w~niniejszej pracy pochodzą z~tego przebiegu. Rzuty poziomów
z~góry wykonano dodatkową kamerą ortograficzną kadrowaną na podstawie obrysu ścian.
Kod skryptu dołączono do materiałów pracy (załącznik~A).

Wszystkie kroki scenariusza zakończyły się powodzeniem, a~w~dzienniku edytora z~fazy
odtwarzania nie zarejestrowano żadnego wyjątku. Pojawiło się natomiast czterokrotnie,
przy wczytywaniu menu i~każdego z~trzech poziomów, ostrzeżenie biblioteki Zenject
o~tworzeniu obiektów przed zakończeniem pracy wszystkich instalatorów. Jest ono skutkiem
przyjętego rozwiązania, w~którym kula jest tworzona wewnątrz metody
\kod{InstallBindings} (listing~\ref{lst:ball-installer}). Rozwiązanie działa poprawnie,
ponieważ wszystkie powiązania wymagane przez kulę są rejestrowane wcześniej, lecz
uzależnia poprawność od kolejności instalatorów w~kontekście sceny. Przeniesienie
tworzenia kuli do etapu inicjalizacji kontenera, np. do klasy implementującej interfejs
\klasa{IInitializable}, usunęłoby tę zależność.

\section{Walidacja konfiguracji i procesu budowania}

Pliki konfiguracyjne są sprawdzane automatycznie przed każdym budowaniem aplikacji
(listing~\ref{lst:walidacja}). Walidacja obejmuje wszystkie 25 klas
konfiguracji, a~błąd w~dowolnym pliku przerywa budowanie z~komunikatem wskazującym klasę
i~przyczynę. W~ten sposób wykrywane są brakujące pliki, sekcje i~wartości, a~także
niezgodność nazw pól JSON z~parametrami konstruktorów klas danych. Proces budowania
dodatkowo weryfikuje, czy wszystkie sceny wskazane w~ustawieniach istnieją w~projekcie.

\section{Testy rozgrywki i wprowadzone poprawki}
\label{sec:poprawki}

Najwięcej zmian w~projekcie wynikło z~sesji testowych rozgrywki, w~których oceniano,
czy mechaniki są zrozumiałe i~czy gra nie stawia gracza w~sytuacjach bez wyjścia.
Zaobserwowane problemy i~wprowadzone na ich podstawie zmiany zestawiono
w~tabeli~\ref{tab:poprawki}. Zmiany te wprowadzano w~kolejnych wersjach projektu.

\begin{table}[!htb]
  \centering
  \caption{Problemy zaobserwowane podczas testów rozgrywki i~wprowadzone zmiany; \zrodloWlasne}
  \label{tab:poprawki}
  \begin{tabularx}{\textwidth}{>{\raggedright\arraybackslash}p{5.6cm} X}
    \toprule
    \textbf{Zaobserwowany problem} & \textbf{Wprowadzona zmiana} \\
    \midrule
    Sterowanie przechyłem (trzy tryby do wyboru w~ustawieniach) okazało się
    nieprecyzyjne & usunięto sterowanie przechyłem i~okno ustawień, pozostawiając wyłącznie
    gest przeciągnięcia \\
    Kula zatrzymana za ścianą bywała niewidoczna i~nie dało się jej kliknąć & dodano
    przycisk \emph{Push}, aktywny tylko przy nieruchomej kuli \\
    Wiatr działający wyłącznie na toczącą się kulę był odbierany jako nieprzewidywalny & wiatr
    stał się stałym, łagodnym dryfem działającym również na stojącą kulę; w~interfejsie
    pozostawiono tylko strzałkę kierunku \\
    Kula dociśnięta wiatrem do ściany obracała się w~miejscu & obrót siatki liczony
    z~rzeczywistego przemieszczenia, a~nie z~zadanej prędkości \\
    Kula stojąca wewnątrz pola przyspieszającego lub apteczki nie reagowała na kliknięcie &
    rzutowanie promienia z~pominięciem wyzwalaczy \\
    Ruchome przeszkody mogły zakleszczyć kulę lub ciągnąć ją za sobą & komponent
    \klasa{ObstacleKnocker} odpychający kulę od punktu styku \\
    Poziomy były odbierane jako puste i~mało dynamiczne & dodano odbijacze, obracające się
    belki i~ruchome ściany oraz zwiększono liczbę pól przyspieszających i~apteczek \\
    Apteczka zebrana przy pełnym zdrowiu przepadała & apteczka pozostaje na planszy, jeśli
    leczenie nie było potrzebne \\
    \bottomrule
  \end{tabularx}
\end{table}

\section{Wydajność i zarządzanie pamięcią}

Najbardziej rozbudowany, trzeci poziom zawiera 89~przeciwników, 82~monety oraz
kilkadziesiąt przeszkód i~pól specjalnych (tabela~\ref{tab:poziomy}). Wszystkie te obiekty,
łącznie z~pulą efektów cząsteczkowych, są tworzone raz, podczas wczytywania poziomu,
a~w~trakcie rozgrywki są jedynie aktywowane i~dezaktywowane. Ruch przeszkód i~przeciwników
jest obliczany w~kroku fizyki, więc jego koszt nie rośnie wraz z~liczbą klatek obrazu.

Drugim źródłem obciążenia w~grach na platformie .NET jest odśmiecanie pamięci. W~kodzie
rozgrywki ograniczono alokacje wykonywane co klatkę: animacje korzystają z~UniTask, którego
typ zadania jest strukturą, subskrypcje strumieni są zakładane jednorazowo przy tworzeniu
obiektów, a~właściwości reaktywne powiadamiają obserwatorów tylko przy rzeczywistej zmianie
wartości. Alokacje występują natomiast przy aktualizacji tekstów interfejsu: licznik czasu
tworzy nowy łańcuch znaków w~każdej klatce, a~pozostałe napisy przy zmianie wyświetlanych
danych. Są to krótkie łańcuchy, ale w~przypadku licznika czasu można je wyeliminować,
korzystając z~metody \kod{SetText} komponentu TextMeshPro, która przyjmuje wartości liczbowe
bez tworzenia pośrednich napisów. Subskrypcje są zwalniane razem z~obiektami, które je
utworzyły; pośrednio potwierdza to scenariusz automatyczny, w~którym kolejne wczytania scen
nie powodowały wyjątków wynikających z~odwołań do zniszczonych obiektów.

W~ramach pracy nie wykonano systematycznych pomiarów liczby klatek na sekundę. Pomiary
takie, wykonane narzędziem Unity Profiler na komputerach o~różnej wydajności, stanowią
naturalny element dalszych prac.
